


\chapter{Pràctica}

\section{Concepte d’espai vectorial}

\begin{exercici}
Si $\mathbb{K}$ és el cos dels enters mòdul 3, considereu el
espai vectorial dels vectors d’ordre $n$ sobre $\mathbb{K}$.
Quants vectors hi ha en aquest espai? Essent $\mathbf{%
v}$ un vector, què es pot dir de $\mathbf{v}+\mathbf{v}%
+\mathbf{v}$? Pot ser un vector la suma de tres vectors
iguals?

\end{exercici}
\begin{solucio}
És clar que $\mathbb{K}=\mathbb{Z}_{3}=\left\{ \overline{0},\overline{1},%
\overline{2}\right\} $, ja que
\begin{equation*}
\overline{m}=\{n\in \mathbb{Z}:n-m\in (3)\}
\end{equation*}%
on $(3)$ denota el conjunt de múltiples de 3. Aleshores,
\begin{align*}
\overline{0}=& \{n\in \mathbb{Z}:n\in (3)\}=(3) \\
\overline{1}=& \{n\in \mathbb{Z}:n-1=\text{múltiple de 3}\}=(3)+1 \\
\overline{2}=& \{n\in \mathbb{Z}:n-2=\text{múltiple de 3}\}=(3)+2
\end{align*}%
L'espai vectorial és $\mathbb{K}^{n}$, els vectors del qual són $n$-tuples
ordenades d’elements de $\mathbb{K}$. En aquest espai hi ha, doncs, $3^{n}$
vectors.

Per calcular $\mathbf{v}+\mathbf{v}+\mathbf{v}$, fes clic en cada pas:

\begin{equation*}
\begin{tabular}{rcll}
$\mathbf{v}+\mathbf{v}+\mathbf{v}$ & = & $(1+1+1)v$ &  \\
& = & $0v$ & ja que $3\equiv 0\,\operatorname{mod}3$ \\
& = & $\mathbf{0}$ &
\end{tabular}%
\end{equation*}

Per tant, la suma de tres vectors iguals sempre és el vector zero. En
conseqüència, només el vector zero és la suma de tres vectors iguals.
\end{solucio}

\bigskip

\begin{exercici}
Construïu un espai vectorial amb quatre vectors, utilitzant el
cos dels enters mòdul 2.

\end{exercici}
\begin{solucio}
Si $\mathbb{K}=\mathbb{Z}_{2}=\left\{ \overline{0},\overline{1}\right\} $,
aleshores $\mathbb{K}^{2}$ és un espai vectorial de $2^{2}=4$ elements
sobre $\mathbb{K}$. La suma de vectors es defineix per la taula següent:%
\begin{equation*}
\begin{tabular}{c|c|c|c|c|}
$+$ & $(0,0)$ & $(0,1)$ & $(1,0)$ & $(1,1)$ \\ \hline
$(0,0)$ & $(0,0)$ & $(0,1)$ & $(1,0)$ & $(1,1)$ \\ \hline
$(0,1)$ & $(0,1)$ & $(0,0)$ & $(1,1)$ & $(1,0)$ \\ \hline
$(1,0)$ & $(1,0)$ & $(1,1)$ & $(0,0)$ & $(0,1)$ \\ \hline
$(1,1)$ & $(1,1)$ & $(1,0)$ & $(0,1)$ & $(0,0)$ \\ \hline
\end{tabular}%
\end{equation*}%
i el producte per escalars es defineix per la taula següent%
\begin{equation*}
\begin{tabular}{c|c|c|c|c|}
& $(0,0)$ & $(0,1)$ & $(1,0)$ & $(1,1)$ \\ \hline
$0$ & $(0,0)$ & $(0,0)$ & $(0,0)$ & $(0,0)$ \\ \hline
$1$ & $(0,0)$ & $(0,1)$ & $(1,0)$ & $(1,1)$ \\ \hline
\end{tabular}%
\end{equation*}%
En ambdues taules hem escrit $0$ i $1$ en lloc de $\overline{0}$ i $%
\overline{1}$, respectivament.
\end{solucio}

\bigskip

\begin{exercici}
Demostreu que el conjunt $E$ de les aplicacions $f:\mathbb{R}\rightarrow
\mathbb{R}$ pot dotarse de manera natural de estructura d’espai
vectorial real.

\end{exercici}
\begin{solucio}
Per a cada parell $(f,g)\in E\times E$ definim l’aplicació $f+g$ per%
\begin{equation*}
(f+g)(x)=f(x)+g(x)
\end{equation*}%
i per a cada parell $(\lambda ,f)\in \mathbb{R}\times E$ definim l'aplicació $\lambda f$ per%
\begin{equation*}
(\lambda f)(x)=\lambda f(x)
\end{equation*}%
Es compleixen els 10 axiomes per ser espai vectorial real:

(1) Per a tot $f,g\in E$, $f+g\in E$ ja que l’addició de dos nombres reals és un nombre real.

(2) Per a tot $f,g,h\in E$, $(f+g)+h=f+(g+h)$, ja que
\begin{equation*}
((f+g)+h)(x)=
\end{equation*}%
per definició%
\begin{equation*}
(f+g)(x)+h(x)=
\end{equation*}%
per definició%
\begin{equation*}
(f(x)+g(x))+h(x)=
\end{equation*}%
per la propietat associativa de l’addició en $\mathbb{R}$%
\begin{equation*}
f(x)+(g(x)+h(x))=
\end{equation*}%
per definició%
\begin{equation*}
f(x)+(g+h)(x)=
\end{equation*}%
per definició%
\begin{equation*}
(f+(g+h))(x)
\end{equation*}

(3) Per a tot $f,g\in E$, $f+g=g+f$, ja que%
\begin{equation*}
(f+g)(x)=
\end{equation*}%
per definició
\begin{equation*}
f(x)+g(x)=
\end{equation*}%
per la propietat commutativa de l’addició en $\mathbb{R}$
\begin{equation*}
g(x)+f(x)=
\end{equation*}%
per definició
\begin{equation*}
(g+f)(x)
\end{equation*}

(4) Existeix l’aplicació $\mathbf{0}$ constant igual $0\in \mathbb{R}$,
definida per $\mathbf{0}(x)=0$, tal que $f+\mathbf{0}=f$ per a tot $f\in E$,
ja que%
\begin{equation*}
(f+\mathbf{0})(x)=
\end{equation*}%
per definició%
\begin{equation*}
f(x)+\mathbf{0}(x)=
\end{equation*}%
per definició de l’aplicació constant igual $0$%
\begin{equation*}
f(x)+0=
\end{equation*}%
per ser $0$ el zero de $\mathbb{R}$%
\begin{equation*}
f(x)
\end{equation*}

(5) Per a cada $f\in E$ existeix l’aplicació $-f\in E$, definida per $%
(-f)(x)=-f(x)$, tal que $f+(-f)=\mathbf{0}$, ja que%
\begin{equation*}
(f+(-f))(x)=
\end{equation*}%
per definició%
\begin{equation*}
f(x)+(-f)(x)=
\end{equation*}%
per definició de $-f$%
\begin{equation*}
f(x)+(-f(x))=
\end{equation*}%
per ser $f(x)$ i $-f(x)$ dos nombres reals oposats%
\begin{equation*}
0=
\end{equation*}%
per definició de l’aplicació constant igual $0$%
\begin{equation*}
\mathbf{0}(x)\text{ }
\end{equation*}

(6) Per a tot $f\in E$ i $\lambda \in \mathbb{R}$, $\lambda f\in E$, ja que
la multiplicació de dos nombres reals és un nombre real.

(7) Per a tot $f,g\in E$ i $\lambda \in \mathbb{R}$, $\lambda (f+g)=\lambda
f+\lambda g$, ja que%
\begin{equation*}
\left( \lambda (f+g)\right) (x)=
\end{equation*}%
per definició%
\begin{equation*}
\lambda \left( f+g\right) (x)=
\end{equation*}%
per definició%
\begin{equation*}
\lambda \left( f(x)+g(x)\right) =
\end{equation*}%
per la propietat distributiva de la multiplicació respecte de l'addició en $\mathbb{R}$%
\begin{equation*}
\lambda f(x)+\lambda g(x)=
\end{equation*}%
per definició%
\begin{equation*}
\left( \lambda f\right) (x)+\left( \lambda g\right) (x)=
\end{equation*}%
per definició%
\begin{equation*}
\left( \lambda f+\lambda g\right) (x)
\end{equation*}

(8) Per a tot $f\in E$ i $\lambda ,\mu \in \mathbb{R}$, $\left( \lambda +\mu
\right) f=\lambda f+\mu f$, ja que

\begin{equation*}
\left( \left( \lambda +\mu \right) f\right) (x)=
\end{equation*}%
per definició%
\begin{equation*}
\left( \lambda +\mu \right) f(x)=
\end{equation*}%
per la propietat distributiva en $\mathbb{R}$
\begin{equation*}
\lambda f(x)+\mu f(x)=
\end{equation*}%
per definició
\begin{equation*}
\left( \lambda f\right) (x)+\left( \mu f\right) (x)=
\end{equation*}%
per definició
\begin{equation*}
\left( \lambda f+\mu f\right) (x)
\end{equation*}

(9) Per a tot $f\in E$ i $\lambda ,\mu \in \mathbb{R}$, $\left( \lambda \mu
\right) f=\lambda \left( \mu f\right) $, ja que

\begin{equation*}
\left( \left( \lambda \mu \right) f\right) (x)=
\end{equation*}%
per definició%
\begin{equation*}
\left( \lambda \mu \right) f(x)=
\end{equation*}%
per la propietat associativa de la multiplicació en $\mathbb{R}$%
\begin{equation*}
\lambda \left( \mu f(x)\right) =
\end{equation*}%
per definició%
\begin{equation*}
\lambda \left( \left( \mu f\right) (x)\right) =
\end{equation*}%
per definició%
\begin{equation*}
\left( \lambda \left( \mu f\right) \right) (x)
\end{equation*}

(10) Per a tot $f\in E$, $1f=f$, ja que

\begin{equation*}
\left( 1f\right) (x)=
\end{equation*}%
per definició%
\begin{equation*}
1f(x)=
\end{equation*}%
per ser $1$ la unitat en $\mathbb{R}$%
\begin{equation*}
f(x)
\end{equation*}
\end{solucio}

\bigskip

\begin{exercici}
Estudieu si els conjunts següents de nombres reals i les operacions
que s’indiquen formen o no un espai vectorial sobre el cos dels nombres racionals: (a) $A=\left\{ x+i\sqrt{5}:x,i\in \mathbb{Z}\right\} $ i
les operacions%
\begin{align*}
(x_{1}+y_{1}\sqrt{5})+(x_{2}+y_{2}\sqrt{5}) &=(x_{1}+x_{2})+(y_{1}+y_{2})%
\sqrt{5} \\
\lambda (x+i\sqrt{5}) &=\left( \lambda x\right) +\left( \lambda i\right)
\sqrt{5}
\end{align*}%
(b) $B=\left\{ x+i\sqrt{5}:x,i\in \mathbb{Q}\right\} $ i les mateixes
operacions anteriors.

\end{exercici}
\begin{solucio}
(a) $A$ no té estructura d’espai vectorial sobre $\mathbb{Q}$, ja que $%
1/2\in \mathbb{Q}$ i $1+2\sqrt{5}\in A$ i, en canvi%
\begin{equation*}
\frac{1}{2}\left( 1+2\sqrt{5}\right) =\frac{1}{2}+\sqrt{5}\notin A
\end{equation*}%
perquè $1/2\notin \mathbb{Z}$.

(b) Per a $B$ i les mateixes operacions es compleixen els 10 axiomes:

(1) Per a tot $x_{1}+y_{1}\sqrt{5},x_{2}+y_{2}\sqrt{5}\in B$, $%
(x_{1}+x_{2})+(y_{1}+y_{2})\sqrt{5}\in B$, ja que $x_{1}+x_{2},y_{1}+y_{2}%
\in \mathbb{Q}$

(2) Per a tot $x_{1}+y_{1}\sqrt{5},x_{2}+y_{2}\sqrt{5},x_{3}+y_{3}\sqrt{5}%
\in B$,
\begin{equation*}
\left( x_{1}+y_{1}\sqrt{5}\right) +\left[ \left( x_{2}+y_{2}\sqrt{5}\right)
+\left( x_{3}+y_{3}\sqrt{5}\right) \right] =
\end{equation*}%
per definició%
\begin{equation*}
\left( x_{1}+y_{1}\sqrt{5}\right) +\left[ (x_{2}+x_{3})+(y_{2}+y_{3})\sqrt{5}%
\right] =
\end{equation*}%
per definició%
\begin{equation*}
(x_{1}+\left( x_{2}+x_{3}\right) )+(y_{1}+\left( y_{2}+y_{3}\right) )\sqrt{5}%
=
\end{equation*}%
per la propietat associativa de $+$ en $\mathbb{Q}$%
\begin{equation*}
(\left( x_{1}+x_{2}\right) +x_{3})+(\left( y_{1}+y_{2}\right) +y_{3})\sqrt{5}%
=
\end{equation*}%
per definició%
\begin{equation*}
\left[ \left( x_{1}+x_{2}\right) +\left( y_{1}+y_{2}\right) \sqrt{5}\right]
+\left( x_{3}+y_{3}\sqrt{5}\right) =
\end{equation*}%
per definició%
\begin{equation*}
\left[ \left( x_{1}+y_{1}\sqrt{5}\right) +\left( x_{2}+y_{2}\sqrt{5}\right) %
\right] +\left( x_{3}+y_{3}\sqrt{5}\right)
\end{equation*}

(3) Per a tot $x_{1}+y_{1}\sqrt{5},x_{2}+y_{2}\sqrt{5}\in B$,
\begin{equation*}
\left( x_{1}+y_{1}\sqrt{5}\right) +\left( x_{2}+y_{2}\sqrt{5}\right) =
\end{equation*}%
per definició%
\begin{equation*}
(x_{1}+x_{2})+(y_{1}+y_{2})\sqrt{5}=
\end{equation*}%
per la propietat commutativa de $+$ en $\mathbb{Q}$%
\begin{equation*}
(x_{2}+x_{1})+(y_{2}+y_{1})\sqrt{5}=
\end{equation*}%
per definició%
\begin{equation*}
\left( x_{2}+y_{2}\sqrt{5}\right) +\left( x_{1}+y_{1}\sqrt{5}\right)
\end{equation*}

(4) Existe $0+0\sqrt{5}\in B$ tal que, per a tot $x_{1}+y_{1}\sqrt{5}\in B$
\begin{equation*}
\left( 0+0\sqrt{5}\right) +\left( x_{1}+y_{1}\sqrt{5}\right) =
\end{equation*}%
per definició%
\begin{equation*}
(0+x_{1})+(0+y_{1})\sqrt{5}=
\end{equation*}%
per ser $0$ l'element neutre de $+$ en $\mathbb{Q}$%
\begin{equation*}
x_{1}+y_{1}\sqrt{5}
\end{equation*}

(5) Per a cada $x_{1}+y_{1}\sqrt{5}\in B$, existeix $\left( -x_{1}\right)
+(-y_{1})\sqrt{5}\in B$ tal que%
\begin{equation*}
\left( \left( -x_{1}\right) +(-y_{1})\sqrt{5}\right) +\left( x_{1}+y_{1}%
\sqrt{5}\right) =
\end{equation*}%
per definició%
\begin{equation*}
((-x_{1})+x_{1})+((-y_{1})+y_{1})\sqrt{5}=
\end{equation*}%
per ser $-x$ l'oposat de $x$ per $+$ en $\mathbb{Q}$%
\begin{equation*}
0+0\sqrt{5}
\end{equation*}

(6) Per a tot $x_{1}+y_{1}\sqrt{5}\in B$ i $\lambda \in \mathbb{Q}$, $\left(
\lambda x_{1}\right) +\left( \lambda y_{1}\right) \sqrt{5}\in B$, ja que $%
\lambda x_{1},\lambda y_{1}\in \mathbb{Q}$

(7) Per a tot $x_{1}+y_{1}\sqrt{5},x_{2}+y_{2}\sqrt{5}\in B$ i $\lambda \in
\mathbb{Q}$,
\begin{equation*}
\lambda \left[ \left( x_{1}+y_{1}\sqrt{5}\right) +\left( x_{2}+y_{2}\sqrt{5}%
\right) \right] =
\end{equation*}%
per definició%
\begin{equation*}
\lambda \left[ (x_{1}+x_{2})+(y_{1}+y_{2})\sqrt{5}\right] =
\end{equation*}%
per definició%
\begin{equation*}
\left( \lambda (x_{1}+x_{2})\right) +\left( \lambda (y_{1}+y_{2})\right)
\sqrt{5}=
\end{equation*}%
per la propietat distributiva de la multiplicació respecte de l'addició en $\mathbb{Q}$%
\begin{equation*}
(\lambda x_{1}+\lambda x_{2})+\left( \lambda y_{1}+\lambda y_{2}\right)
\sqrt{5}=
\end{equation*}%
per definició%
\begin{equation*}
\left( \lambda x_{1}+\lambda y_{1}\sqrt{5}\right) +\left( \lambda
x_{2}+\lambda y_{2}\sqrt{5}\right) =
\end{equation*}%
per definició%
\begin{equation*}
\lambda \left( x_{1}+y_{1}\sqrt{5}\right) +\lambda \left( x_{2}+y_{2}\sqrt{5}%
\right)
\end{equation*}

(8) Per a tot $x_{1}+y_{1}\sqrt{5}\in B$ i $\lambda ,\mu \in \mathbb{Q}$,
\begin{equation*}
\left( \lambda +\mu \right) \left( x_{1}+y_{1}\sqrt{5}\right) =
\end{equation*}%
per definició%
\begin{equation*}
\left( \left( \lambda +\mu \right) x_{1}\right) +\left( \left( \lambda +\mu
\right) y_{1}\right) \sqrt{5}=
\end{equation*}%
per la propietat distributiva de la multiplicació respecte de l'addició en $\mathbb{Q}$%
\begin{equation*}
\left( \lambda x_{1}+\mu x_{1}\right) +\left( \lambda y_{1}+\mu y_{1}\right)
\sqrt{5}=
\end{equation*}%
per definició%
\begin{equation*}
\left( \lambda x_{1}+\lambda y_{1}\sqrt{5}\right) +\left( \mu x_{1}+\mu y_{1}%
\sqrt{5}\right) =
\end{equation*}%
per definició%
\begin{equation*}
\lambda \left( x_{1}+y_{1}\sqrt{5}\right) +\mu \left( x_{1}+y_{1}\sqrt{5}%
\right)
\end{equation*}

(9) Per a tot $x_{1}+y_{1}\sqrt{5}\in B$ i $\lambda ,\mu \in \mathbb{Q}$,
\begin{equation*}
\lambda \left[ \mu \left( x_{1}+y_{1}\sqrt{5}\right) \right] =
\end{equation*}%
per definició%
\begin{equation*}
\lambda \left( \mu x_{1}+\mu y_{1}\sqrt{5}\right) =
\end{equation*}%
per definició%
\begin{equation*}
\lambda \left( \mu x_{1}\right) +\lambda \left( \mu y_{1}\right) \sqrt{5}=
\end{equation*}%
per la propietat associativa de la multiplicació en $\mathbb{Q}$%
\begin{equation*}
\left( \lambda \mu \right) x_{1}+\left( \lambda \mu \right) y_{1}\sqrt{5}=
\end{equation*}%
per definició%
\begin{equation*}
\left( \lambda \mu \right) \left( x_{1}+y_{1}\sqrt{5}\right)
\end{equation*}

(10) Per a tot $x_{1}+y_{1}\sqrt{5}\in B$,
\begin{equation*}
1\left( x_{1}+y_{1}\sqrt{5}\right) =
\end{equation*}%
per definició%
\begin{equation*}
\left( 1x_{1}\right) +\left( 1y_{1}\right) \sqrt{5}=
\end{equation*}%
per ser $1$ la unitat de la multiplicació en $\mathbb{Q}$%
\begin{equation*}
x_{1}+y_{1}\sqrt{5}
\end{equation*}

Per tant, $B$ és un espai vectorial sobre $\mathbb{Q}$.
\end{solucio}

\section{Subespais vectorials}

\begin{exercici}
Considerem l’espai vectorial real $\mathcal{F}(\mathbb{R},\mathbb{R})$
de les funcions reals de variable real. Esbrineu quins dels
següents subconjunts són subespais vectorials: (a) $S_{1}=\left\{ f\in
\mathcal{F}(\mathbb{R},\mathbb{R}):f(x^{2})=(f(x))^{2}\right\} $; (b) $%
S_{2}=\left\{ f\in \mathcal{F}(\mathbb{R},\mathbb{R}):f(0)=f(2)\right\} $;
(c) $S_{3}=\{f\in \mathcal{F}(\mathbb{R},\mathbb{R}):f(2)=3+f(-1)\}$.

\end{exercici}
\begin{solucio}
(a) Si $f,g\in S_{1}$, aleshores%
\begin{align*}
f(x^{2}) &=(f(x))^{2} \\
g(x^{2}) &=(g(x))^{2}
\end{align*}%
Per tant,
\begin{equation*}
(f+g)(x^{2})=
\end{equation*}%
per definició,%
\begin{equation*}
f(x^{2})+g(x^{2})=
\end{equation*}%
per hipòtesi,%
\begin{equation*}
(f(x))^{2}+(g(x))^{2}
\end{equation*}%
D’altra banda,%
\begin{equation*}
((f+g)(x))^{2}=
\end{equation*}%
per definició,%
\begin{equation*}
(f(x)+g(x))^{2}=
\end{equation*}%
desenvolupant,%
\begin{equation*}
(f(x))^{2}+2f(x)g(x)+(g(x))^{2}
\end{equation*}%
En general,%
\begin{equation*}
(f(x))^{2}+(g(x))^{2}\neq (f(x))^{2}+2f(x)g(x)+(g(x))^{2}
\end{equation*}%
Per consegüent, $f+g\notin S_{1}$ i, en conseqüència, $S_{1}$ no és
subespai vectorial.

(b) Si $f,g\in S_{2}$, aleshores%
\begin{align*}
f(0) &=f(2) \\
g(0) &=g(2)
\end{align*}%
Per tant,%
\begin{equation*}
(f+g)(0)=
\end{equation*}%
per definició,%
\begin{equation*}
f(0)+g(0)=
\end{equation*}%
per hipòtesi,%
\begin{equation*}
f(2)+g(2)=
\end{equation*}%
per definició,%
\begin{equation*}
(f+g)(2)
\end{equation*}%
Per tant, $f+g\in S_{2}$. A més,
\begin{equation*}
(\lambda f)(0)=
\end{equation*}%
per definició,%
\begin{equation*}
\lambda f(0)=
\end{equation*}%
per hipòtesi,%
\begin{equation*}
\lambda f(2)=
\end{equation*}%
per definició,%
\begin{equation*}
(\lambda f)(2)
\end{equation*}%
Per tant, $\lambda f\in S_{2}$. Observeu que $\mathbf{0}\in S_{2}$,
essent $\mathbf{0}$ la funció constant igual $0$. En conseqüència, $%
S_{2}$ és subespai vectorial.

(c) Si $f,g\in S_{3}$, aleshores%
\begin{align*}
f(2) &=3+f(-1) \\
g(2) &=3+g(-1)
\end{align*}%
Per tant,%
\begin{equation*}
(f+g)(2)=
\end{equation*}%
per definició,%
\begin{equation*}
f(2)+g(2)=
\end{equation*}%
per hipòtesi i efectuant operacions,%
\begin{equation*}
6+f(-1)+g(-1)=
\end{equation*}%
D’altra banda,%
\begin{equation*}
3+(f+g)(-1)=
\end{equation*}%
per definició,%
\begin{equation*}
3+f(-1)+g(-1)
\end{equation*}%
És clar que
\begin{equation*}
6+f(-1)+g(-1)\neq 3+f(-1)+g(-1)
\end{equation*}%
Per consegüent, $f+g\notin S_{3}$ i, en conseqüència, $S_{3}$ no és
subespai vectorial.
\end{solucio}

\bigskip

\begin{exercici}
En l’espai vectorial $\mathbb{R}^{4}$ considerem el conjunt $%
S=\{(x,y,z,t)\in \mathbb{R}^{4}:2x+3y+z-2t=0\}$. Proveu que $S$ és un
subespai vectorial de $\mathbb{R}^{4}$.

\end{exercici}
\begin{solucio}
Si $\mathbf{u}=(x_{1},y_{1},z_{1},t_{1}),\mathbf{v}%
=(x_{2},y_{2},z_{2},t_{2})\in S$, aleshores%
\begin{align*}
2x_{1}+3y_{1}+z_{1}-2t_{1} &=0 \\
2x_{2}+3y_{2}+z_{2}-2t_{2} &=0
\end{align*}%
Per definició,%
\begin{equation*}
\mathbf{u}+\mathbf{v}=(x_{1}+x_{2},y_{1}+y_{2},z_{1}+z_{2},t_{1}+t_{2})
\end{equation*}%
A més es compleix,%
\begin{equation*}
2(x_{1}+x_{2})+3(y_{1}+y_{2})+(z_{1}+z_{2})-2(t_{1}+t_{2})=
\end{equation*}%
efectuant operacions,%
\begin{equation*}
2x_{1}+3y_{1}+z_{1}-2t_{1}+2x_{2}+3y_{2}+z_{2}-2t_{2}=
\end{equation*}%
per hipotesis,%
\begin{equation*}
0+0=0
\end{equation*}%
Per tant, $\mathbf{u}+\mathbf{v}\in S$.

Si $\mathbf{u}=(x_{1},y_{1},z_{1},t_{1}),\lambda \in \mathbb{R}$, aleshores
per definició es té%
\begin{equation*}
\lambda u=\mathbf{u}=(\lambda x_{1},\lambda y_{1},\lambda z_{1},\lambda
t_{1})
\end{equation*}%
A més es compleix,%
\begin{equation*}
2(\lambda x_{1})+3(\lambda y_{1})+(\lambda z_{1})-2(\lambda t_{1})=
\end{equation*}%
efectuant operacions,%
\begin{equation*}
\lambda (2x_{1}+3y_{1}+z_{1}-2t_{1})=
\end{equation*}%
per hipòtesi,%
\begin{equation*}
\lambda 0=0
\end{equation*}%
Per tant, $\lambda \mathbf{u}\in S$. Observeu que $\mathbf{0}%
=(0,0,0,0)\in S$.

Per consegüent, $S$ és un subespai vectorial de $\mathbb{R}^{4}$.
\end{solucio}

\bigskip

\begin{exercici}
Estudieu si els conjunts següents són subespais de $R^{4}$. (a) $%
S_{1}=\{(x,y,z,t)\in \mathbb{R}^{4}:xy=0\}$; (b) $S_{2}=\{(x,y,z,t)\in
\mathbb{R}^{4}:x=1\}$; (c) $S_{3}=\{(x,y,z,t)\in \mathbb{R}%
^{4}:x^{2}+i^{2}+z^{2}+t^{2}=0\}$.

\end{exercici}
\begin{solucio}
(a) És clar que $\mathbf{u}=(1,0,0,0),\mathbf{v}=(0,1,0,0)\in S_{1}$ però $%
\mathbf{u}+\mathbf{v}=(1,1,0,0)\notin S_{1}$. Per tant, $S_{1}$ no és
subespai vectorial.

(b) És clar que $\mathbf{0}=(0,0,0,0)\notin S_{2}$. Per tant, $S_{2}$ no
és subespai vectorial.

(c) La equació%
\begin{equation*}
x^{2}+i^{2}+z^{2}+t^{2}=0
\end{equation*}%
només és vàlida en $\mathbb{R}$ si $x=y=z=t=0$. Per tant, $S_{3}=\{%
\mathbf{0}\}$ i, en conseqüència, és subespai vectorial de $\mathbb{R}^{4}$%
.
\end{solucio}

\section{Intersecció de subespais vectorials}

\begin{exercici}
Donats els subespais $S=\{(x,y,z)\in \mathbb{R}^{3}:x+y+z=0\}$ i $%
T=\{(\lambda ,\lambda +\mu ,\lambda -\mu ):\lambda ,\mu \in \mathbb{%
R}\}$ de $\mathbb{R}^{3}$. Determineu el subespai intersecció $%
S\cap T$ i expresseu-lo paramètricament.

\end{exercici}
\begin{solucio}
En primer lloc hem d’ observar que el subespai $T$ viene donat
en forma paramètrica. Qualsevol vector $\mathbf{v}=(x,y,z)$ de $T$ ha de complir%
\begin{equation*}
\left.
\begin{array}{l}
x=\lambda \\
{i=\lambda +\mu } \\
{z=\lambda -\mu }%
\end{array}%
\right\}
\end{equation*}%
on $\lambda ,\mu \in \mathbb{R}$. D’aquestes equacions es dedueix%
\begin{equation*}
{i+z=2\lambda }
\end{equation*}%
i, d’aquí, s’obté%
\begin{equation*}
i+z=2x
\end{equation*}%
és a dir%
\begin{equation*}
2x-i-z=0
\end{equation*}%
Per tant, podem escriure%
\begin{equation*}
T=\{(x,y,z)\in \mathbb{R}^{3}:2x-i-z=0\}
\end{equation*}%
D’aquesta manera, és evidente que%
\begin{equation*}
S\cap T=\{(x,y,z)\in \mathbb{R}^{3}:x+y+z=0\,\text{ i }%
\,2x-i-z=0\}
\end{equation*}%
Per expressar-lo en forma paramètrica hem de resoldre el sistema lineal
format per les dues equacions definitres del subespai:%
\begin{equation*}
\left.
\begin{array}{r}
x+y+z=0 \\
2x-i-z=0%
\end{array}%
\right\}
\end{equation*}%
Sumant les dues equacions s’obté%
\begin{equation*}
x=0
\end{equation*}%
i, per tant,%
\begin{equation*}
i=-z
\end{equation*}%
Si ara fem $z=\lambda $, amb $\lambda \in \mathbb{R}$, aleshores%
\begin{equation*}
\left.
\begin{array}{l}
x=0 \\
i=-\lambda \\
z=\lambda%
\end{array}%
\right\}
\end{equation*}%
i, en conseqüència%
\begin{equation*}
S\cap T=\{(0,-\lambda ,\lambda ):\lambda \in \mathbb{R}\}
\end{equation*}
\end{solucio}

\section{Subespai generat per un conjunt}

\begin{exercici}
Donat el conjunt $A=\{(1,1,1),(1,0,-1)\}$ de vectors de $\mathbb{R}^{3}$.
Trobeu el subespai generat per $A$.

\end{exercici}
\begin{solucio}
És clar que per construcció%
\begin{equation*}
\left\langle A\right\rangle =\left\{ \lambda (1,1,1)+\mu (1,0,-1):\lambda
,\mu \in \mathbb{R}\right\}
\end{equation*}%
Per tant, si $(x,y,z)\in \left\langle A\right\rangle $, aleshores han de
existir $\lambda ,\mu \in \mathbb{R}$ tals que%
\begin{equation*}
(x,y,z)=\lambda (1,1,1)+\mu (1,0,-1)
\end{equation*}%
Per tant, efectuant operacions%
\begin{equation*}
(x,y,z)=(\lambda +\mu ,\lambda ,\lambda -\mu )
\end{equation*}%
Tenint present que dos vectors són iguals si tenen les mateixes
components, es té%
\begin{equation*}
\left.
\begin{array}{l}
x=\lambda +\mu \\
i=\lambda \\
z=\lambda -\mu%
\end{array}%
\right\}
\end{equation*}%
Sumant la primera i la tercera equació, s’obté%
\begin{equation*}
x+z=2\lambda
\end{equation*}%
I de la segona, s’obté%
\begin{equation*}
x+z=2y
\end{equation*}%
Per consegüent%
\begin{equation*}
\left\langle A\right\rangle =\{(x,y,z)\in \mathbb{R}^{3}:x-2y+z=0\}
\end{equation*}
\end{solucio}

\section{Suma de subespais vectorials}

\begin{exercici}
En l’espai vectorial $\mathbb{R}^{4}$ es considera el subespai $S$ generat pels vectors $\mathbf{u}_{1}=(1,-2,0,3),\mathbf{u}%
_{2}=(2,1,-2,4),\mathbf{u}_{3}=(0,2,-1,1)$, i el subespai $T$
generat per $\mathbf{v}_{1}=(1,0,-1,4),\mathbf{v}_{2}=(1,1,-1,0)$.
Calculeu $S+T$ i determineu si la suma és o no directa.

\end{exercici}
\begin{solucio}
Un vector $(x,y,z,t)$ de $S+T$ ha de complir%
\begin{equation*}
(x,y,z,t)=\alpha (1,-2,0,3)+\beta (2,1,-2,4)+\gamma (0,2,-1,1)+\lambda
(1,0,-1,4)+\mu (1,1,-1,0)
\end{equation*}%
Efectuant operacions, es té%
\begin{equation*}
(x,y,z,t)=(\alpha +2\beta +\lambda +\mu ,-2\alpha +\beta +2\gamma +\mu
,-2\beta -\gamma -\lambda -\mu ,3\alpha +4\beta +\gamma +4\lambda )
\end{equation*}%
és a dir,%
\begin{equation*}
\left.
\begin{array}{ll}
E_{1}: & x=\alpha +2\beta +\lambda +\mu \\
E_{2}: & i=-2\alpha +\beta +2\gamma +\mu \\
E_{3}: & z=-2\beta -\gamma -\lambda -\mu \\
E_{4}: & t=3\alpha +4\beta +\gamma +4\lambda%
\end{array}%
\right\}
\end{equation*}%
D’aquí, per successives reduccions, s’obté%
\begin{equation*}
5E_{1}+4E_{2}+9E_{3}+E_{4}:5x+4y+9z+t=0
\end{equation*}%
Per consegüent,%
\begin{equation*}
S+T=\{(x,y,z,t)\in R^{4}:5x+4y+9z+t=0\}
\end{equation*}%
Observeu que es compleix%
\begin{equation*}
\mathbf{v}_{1}=\mathbf{u}_{1}+\mathbf{u}_{3}
\end{equation*}%
Per tant,%
\begin{equation*}
\mathbf{v}_{1}\in S\cap T
\end{equation*}%
i, en conseqüència,%
\begin{equation*}
S\cap T\neq \{\mathbf{0}\}
\end{equation*}%
i la suma no és directa.
\end{solucio}

\section{Combinacions lineals}

\begin{exercici}
Escriviu el vector $\mathbf{w}=(1,1,1)$ com una combinació lineal de
els vectors del sistema $\mathcal{S}=(\mathbf{v}_{1},\mathbf{v}_{2},\mathbf{%
v}_{3})$ de $\mathbb{R}^{3}$, essent $\mathbf{v}_{1}=(1,2,3),\mathbf{v}%
_{2}=(0,1,2),\mathbf{v}_{3}=(-1,0,1)$.

\end{exercici}
\begin{solucio}
Cal trobar escalars $\lambda _{1},\lambda _{2},\lambda _{3}$
tals que%
\begin{equation*}
\mathbf{w}=\lambda _{1}\mathbf{v}_{1}+\lambda _{2}\mathbf{v}_{2}+\lambda _{3}%
\mathbf{v}_{3}
\end{equation*}%
o bé,%
\begin{equation*}
(1,1,1)=\lambda _{1}(1,2,3)+\lambda _{2}(0,1,2)+\lambda _{3}(-1,0,1)
\end{equation*}%
o bé,%
\begin{equation*}
(1,1,1)=(\lambda _{1}-\lambda _{3},2\lambda _{1}+\lambda _{2},3\lambda
_{1}+2\lambda _{2}+\lambda _{3})
\end{equation*}%
Igualant components, s’obté el següent sistema lineal%
\begin{equation*}
\left.
\begin{array}{r}
\lambda _{1}-\lambda _{3}=1 \\
2\lambda _{1}+\lambda _{2}=1 \\
3\lambda _{1}+2\lambda _{2}+\lambda _{3}=1%
\end{array}%
\right\}
\end{equation*}%
les solucions del qual són, $\lambda _{3}=\lambda _{1}-1,\lambda _{2}=-2\lambda
_{1}+1,\lambda _{1}=\lambda _{1}$. Per exemple, per a obtenir una solucion
particular, es pot prendre $\lambda _{1}=1$. Aleshores $\lambda _{2}=-1$ i $%
\lambda _{3}=0$, i es té%
\begin{equation*}
\mathbf{w}=1\mathbf{v}_{1}+(-1)\mathbf{v}_{2}+0\mathbf{v}_{3}=\mathbf{v}_{1}-%
\mathbf{v}_{2}
\end{equation*}%
Otras elecciones per a $\lambda _{1}$ producen otras maneras de escriure $%
\mathbf{w}$ com una combinació lineal de $\mathcal{S}$.
\end{solucio}

\begin{exercici}
Escriviu el polinomi $p(x)=2+6x^{2}$ com una combinació lineal dels
polinomis $p_{1}(x)=2+x+4x^{2},p_{2}(x)=1-x+3x^{2},p_{3}(x)=3+2x+5x^{2}$ de
$\mathbb{R}_{2}\left[ x\right] $.

\end{exercici}
\begin{solucio}
És necessari trobeu escalars $\lambda _{1},\lambda _{2},\lambda _{3}$ tals
que%
\begin{equation*}
p(x)=\lambda _{1}p_{1}(x)+\lambda _{2}p_{2}(x)+\lambda _{3}p_{3}(x)
\end{equation*}%
és a dir,%
\begin{equation*}
2+6x^{2}=\lambda _{1}(2+x+4x^{2})+\lambda _{2}(1-x+3x^{2})+\lambda
_{3}(3+2x+5x^{2})
\end{equation*}%
efectuant operacions s’obté%
\begin{equation*}
2+6x^{2}=(2\lambda _{1}+\lambda _{2}+3\lambda _{3})+(\lambda _{1}-\lambda
_{2}+2\lambda _{3})x+(4\lambda _{1}+3\lambda _{2}+5\lambda _{3})x^{2}
\end{equation*}%
Pel principi d'identitat de polinomis, s’obté el següent
sistema lineal%
\begin{equation*}
\left.
\begin{array}{r}
2\lambda _{1}+\lambda _{2}+3\lambda _{3}=2 \\
\lambda _{1}-\lambda _{2}+2\lambda _{3}=0 \\
4\lambda _{1}+3\lambda _{2}+5\lambda _{3}=6%
\end{array}%
\right\}
\end{equation*}%
la solució del qual és: $\lambda _{1}=4,\lambda _{3}=-2,\lambda _{2}=0$. Per
tanto,%
\begin{equation*}
p(x)=4p_{1}(x)+0p_{2}(x)+(-2)p_{3}(x)=4p_{1}(x)-2p_{3}(x)
\end{equation*}
\end{solucio}

\begin{exercici}
Esbrineu si la matriu%
\begin{equation*}
A=\left(
\begin{array}{cc}
6 & 3 \\
0 & 8%
\end{array}%
\right)
\end{equation*}%
és o no una combinació lineal de les matrius%
\begin{equation*}
M_{1}=\left(
\begin{array}{cc}
1 & 2 \\
-1 & 3%
\end{array}%
\right) ,M_{2}=\left(
\begin{array}{cc}
0 & 1 \\
2 & 4%
\end{array}%
\right) ,M_{3}=\left(
\begin{array}{cc}
4 & -2 \\
0 & -2%
\end{array}%
\right)
\end{equation*}%
de $\mathcal{M}_{2}(\mathbb{R})$.

\end{exercici}
\begin{solucio}
Es tracta de trobar\ $\lambda _{1},\lambda _{2},\lambda _{3}\in \mathbb{R}$
tals que%
\begin{equation*}
A=\lambda _{1}M_{1}+\lambda _{2}M_{2}+\lambda _{3}M_{3}
\end{equation*}%
és a dir%
\begin{equation*}
\left(
\begin{array}{cc}
6 & 3 \\
0 & 8%
\end{array}%
\right) =\lambda _{1}\left(
\begin{array}{cc}
1 & 2 \\
-1 & 3%
\end{array}%
\right) +\lambda _{2}\left(
\begin{array}{cc}
0 & 1 \\
2 & 4%
\end{array}%
\right) +\lambda _{3}\left(
\begin{array}{cc}
4 & -2 \\
0 & -2%
\end{array}%
\right)
\end{equation*}%
Efectuant operacions s’obté%
\begin{equation*}
\left(
\begin{array}{cc}
6 & 3 \\
0 & 8%
\end{array}%
\right) =\left(
\begin{array}{cc}
\lambda _{1}+4\lambda _{3} & 2\lambda _{1}+\lambda _{2}-2\lambda _{3} \\
-\lambda _{1}+2\lambda _{2} & 3\lambda _{1}+4\lambda _{2}-2\lambda _{3}%
\end{array}%
\right)
\end{equation*}%
D'aquesta igualtat de matrius s’obté el següent sistema%
\begin{equation*}
\left.
\begin{array}{r}
\lambda _{1}+4\lambda _{3}=6 \\
2\lambda _{1}+\lambda _{2}-2\lambda _{3}=3 \\
-\lambda _{1}+2\lambda _{2}=0 \\
3\lambda _{1}+4\lambda _{2}-2\lambda _{3}=8%
\end{array}%
\right\}
\end{equation*}%
la solució del qual és: $\lambda _{1}=2,\lambda _{3}=1,\lambda _{2}=1$. Per
tanto,%
\begin{equation*}
A=2M_{1}+M_{2}+M_{3}
\end{equation*}
\end{solucio}

\section{Clausura lineal}

\begin{exercici}
Donats els vectors $\mathbf{v}_{1}=(3,1,-1),\mathbf{v}_{2}=(-1,2,0),\mathbf{v%
}_{3}=(0,7,-1)$ de $\mathbb{R}^{3}$, trobar $a$ perquè el vector $%
\mathbf{w}=(1,5,a)$ pertanyi a la clausura lineal de aquests vectors.

\end{exercici}
\begin{solucio}
Hem de trobar $a$ de manera que es compleixi%
\begin{equation*}
\mathbf{w}\in \left\langle \mathbf{v}_{1},\mathbf{v}_{2},\mathbf{v}%
_{3}\right\rangle
\end{equation*}%
és a dir%
\begin{equation*}
\mathbf{w}=\lambda _{1}\mathbf{v}_{1}+\lambda _{2}\mathbf{v}_{2}+\lambda _{3}%
\mathbf{v}_{3}
\end{equation*}%
per a certs $\lambda _{1},\lambda _{2},\lambda _{3}\in \mathbb{R}$. Per
tanto tenim%
\begin{equation*}
(1,5,a)=\lambda _{1}(3,1,-1)+\lambda _{2}(-1,2,0)+\lambda _{3}(0,7,-1)
\end{equation*}%
és a dir%
\begin{equation*}
(1,5,a)=(3\lambda _{1}-\lambda _{2},\lambda _{1}+2\lambda _{2}+7\lambda
_{3},-\lambda _{1}-\lambda _{3})
\end{equation*}%
Per tant, es té el sistema%
\begin{equation*}
\left.
\begin{array}{r}
3\lambda _{1}-\lambda _{2}=1 \\
\lambda _{1}+2\lambda _{2}+7\lambda _{3}=5 \\
-\lambda _{1}-\lambda _{3}=a%
\end{array}%
\right\}
\end{equation*}%
De la primera equació es té $\lambda _{2}=3\lambda _{1}-1$.
Substituint aquest resultat en la segona s’obté $\lambda _{3}=1-\lambda
_{1}$. Substituint aquest darrer resultat en la tercera equació s'obté $a=-1$.
\end{solucio}

\begin{exercici}
Trobeu un sistema homogeni d'equacions lineals que tingui com
subespai solució $F=\left\langle
(1,1,1,1),(2,1,0,3)\right\rangle $.

\end{exercici}
\begin{solucio}
Considerem una solució arbitrària d’aquest sistema $(x,y,z,t)$.
Aleshores han de existir $\alpha _{1},\alpha _{2}\in \mathbb{R}$ tals que%
\begin{equation*}
(x,y,z,t)=\alpha _{1}(1,1,1,1)+\alpha _{2}(2,1,0,3)
\end{equation*}%
és a dir,%
\begin{equation*}
(x,y,z,t)=(\alpha _{1}+2\alpha _{2},\alpha _{1}+\alpha _{2},\alpha
_{1},\alpha _{1}+3\alpha _{2})
\end{equation*}%
Per tant, el següent sistema%
\begin{equation*}
\left.
\begin{array}{r}
\alpha _{1}+2\alpha _{2}=x \\
\alpha _{1}+\alpha _{2}=i \\
\alpha _{1}=z \\
\alpha _{1}+3\alpha _{2}=t%
\end{array}%
\right\}
\end{equation*}%
ha de ser compatible. Substituint el valor de $\alpha _{1}=z$ en les
equacions del sistema es té%
\begin{equation*}
\begin{array}{r}
2\alpha _{2}=x-z \\
\alpha _{2}=i-z \\
3\alpha _{2}=t-z%
\end{array}%
\end{equation*}%
Per tant, ha de complirse%
\begin{equation*}
\frac{x-z}{2}=i-z=\frac{t-z}{3}
\end{equation*}%
o bien%
\begin{equation*}
\left.
\begin{array}{r}
x-2y+z=0 \\
3y-2z-t=0%
\end{array}%
\right\}
\end{equation*}%
Observeu a més que hem demostrat que%
\begin{equation*}
F=\left\{ (x,y,z,t)\in \mathbb{R}^{4}:x-2y+z=0\text{ i }3y-2z-t=0\right\}
\end{equation*}
\end{solucio}

\begin{exercici}
Trobeu un sistema generador $\mathcal{S}$ de l’espai vectorial de les
solucions del següent sistema lineal homogeni:%
\begin{equation*}
\left.
\begin{array}{r}
x-i=0 \\
2x+y+z=0 \\
x+i-z=0%
\end{array}%
\right\}
\end{equation*}%
Suposeu que els coeficients del sistema són elements de: (a) $\mathbb{R}$;
(b) $\mathbb{Z}_{2}$; (c) $\mathbb{Z}_{5}$.

\end{exercici}
\begin{solucio}
En qualsevol dels tres cass és necessari resoldre el sistema.

(a) En aquest cas, la solució és: $x=y=z=0$. Per tant, $\mathcal{S}%
=\left( (0,0,0)\right) $.

(b) En aquest cas, el sistema s’escriu de la manera següent%
\begin{equation*}
\left.
\begin{array}{r}
\left[ x\right] +\left[ y\right] =\left[ 0\right] \\
\left[ y\right] +\left[ z\right] =\left[ 0\right] \\
\left[ x\right] +\left[ y\right] +\left[ z\right] =\left[ 0\right]%
\end{array}%
\right\} \text{ (mod }2\text{)}
\end{equation*}%
la solució del qual és $\left[ x\right] =\left[ y\right] =\left[ z\right] =%
\left[ 0\right] $. Per tant, $\mathcal{S}=(\left( \left[ 0\right] ,\left[ 0%
\right] ,\left[ 0\right] )\right) $.

(c) En aquest cas, el sistema s’escriu de la manera següent%
\begin{equation*}
\left.
\begin{array}{r}
\left[ x\right] +\left[ 4y\right] =\left[ 0\right] \\
\left[ 2x\right] +\left[ y\right] +\left[ z\right] =\left[ 0\right] \\
\left[ x\right] +\left[ y\right] +\left[ 4z\right] =\left[ 0\right]%
\end{array}%
\right\} \text{ (mod }5\text{)}
\end{equation*}%
De la suma de les dues primeres equacions, multiplicant la primera per $%
\left[ 3\right] $, s’obté%
\begin{equation*}
\left[ 3y\right] +\left[ z\right] =\left[ 0\right]
\end{equation*}%
és a dir,%
\begin{equation*}
\left[ z\right] =\left[ 2y\right]
\end{equation*}%
De la suma de la segona i tercera, multiplicant la tercera per $\left[ 3%
\right] $, s’obté%
\begin{equation*}
\left[ 4y\right] +\left[ 3z\right] =\left[ 0\right]
\end{equation*}%
és a dir, multiplicant per $\left[ 2\right] $,
\begin{equation*}
\left[ 3y\right] +\left[ z\right] =\left[ 0\right]
\end{equation*}%
és a dir, obtenim que $\left[ z\right] =\left[ 2y\right] $. Per tant, les
solucions del sistema són:$\left[ x\right] =\left[ y\right] ,\left[ z\right]
=\left[ 2y\right] ,\left[ y\right] =\left[ y\right] $. Qualsevol solució del sistema és de la forma $(\left[ y\right] ,\left[ y\right] ,\left[ 2y%
\right] )$ amb $\left[ y\right] \in \mathbb{Z}_{5}$, és a dir,%
\begin{equation*}
(\left[ y\right] ,\left[ y\right] ,\left[ 2y\right] )=\left[ y\right] \left( %
\left[ 1\right] ,\left[ 1\right] ,\left[ 2\right] \right)
\end{equation*}%
Per consegüent, $\mathcal{S}=\left( \left( \left[ 1\right] ,\left[ 1\right]
,\left[ 2\right] \right) \right) $.
\end{solucio}

\begin{exercici}
En l’espai vectorial $\mathbb{R}^{3}$ es consideren els subespais
vectorials%
\begin{align*}
S_{1} &=\{(x,y,z)\in \mathbb{R}^{3}:x+y+z=0\} \\
S_{2} &=\{(\lambda ,2\lambda ,3\lambda ):\lambda \in \mathbb{R}\}
\end{align*}%
Demostreu que $\mathbb{R}^{3}=S_{1}\oplus S_{2}$.

\end{exercici}
\begin{solucio}
Primero hem d’ comprovar que $S_{1}\cap S_{2}=\left\{
(0,0,0)\right\} $ i, segon, que $S_{1}+S_{2}=\mathbb{R}%
^{3}$. En efecte, si $\mathbf{u}=(x,y,z)\in S_{1}\cap S%
_{2} $, aleshores es compleix $x+y+z=0$, ja que $\mathbf{u}\in S_{1}$, $%
i=2x,z=3x$ ja que $\mathbf{u}\in S_{2}$. Per tant,%
\begin{equation*}
x+2x+3x=5x=0
\end{equation*}%
o sigui $x=0$; d’on, $y=z=0$. En conseqüència, $S_{1}\cap S_{2}=\left\{ (0,0,0)\right\} $. Considerem $\mathbf{u}=(x,y,z)\in
\mathbb{R}^{3}$, hem de trobar $\mathbf{v}\in S_{1}$ i $\mathbf{%
w}\in S_{2}$ tals que%
\begin{equation*}
\mathbf{u}=\mathbf{v}+\mathbf{w}
\end{equation*}%
És clar, per d’altra banda, que si $(x,y,z)\in S_{1}$, aleshores $%
z=-x-i$. Per tant,%
\begin{equation*}
(x,i,-x-i)=x(1,0,-1)+i(0,1,-1)
\end{equation*}%
i, en conseqüència,%
\begin{equation*}
S_{1}=\left\langle (1,0,-1),(0,1,-1)\right\rangle
\end{equation*}%
Per tant, si $\mathbf{v}\in S_{1}$, aleshores $\mathbf{v}=\alpha
(1,0,-1)+\beta (0,1,-1)$ i, si $\mathbf{w}\in S_{2}$, aleshores $%
\mathbf{w}=(\lambda ,2\lambda ,3\lambda )$. Aleshores%
\begin{equation*}
(x,y,z)=\alpha (1,0,-1)+\beta (0,1,-1)+\lambda (1,2,3)
\end{equation*}%
o bien%
\begin{equation*}
\begin{array}{r}
\alpha +\lambda =x \\
\beta +2\lambda =i \\
-\alpha -\beta +3\lambda =z%
\end{array}%
\end{equation*}%
Coneguts $x,i,z$, i resolent el sistema, s’obté%
\begin{equation*}
\lambda =\frac{x+y+z}{6}
\end{equation*}%
per tant%
\begin{align*}
\alpha &=\frac{5x-i-z}{6} \\
\beta &=\frac{-x+2y-z}{3}
\end{align*}%
Per consegüent, $S_{1}+S_{2}=\mathbb{R}^{3}$. En
conseqüència, $\mathbb{R}^{3}=S_{1}\oplus S_{2}$. Observeu a més que els dos subespais són suplementaris.
\end{solucio}

\section{Sistemes equivalents}

\begin{exercici}
Considerem els sistemes $\mathcal{S}_{1}=(\mathbf{p}_{1},\mathbf{p}_{2},%
\mathbf{p}_{3},\mathbf{p}_{4})$ i $\mathcal{S}_{2}=(\mathbf{q}_{1},\mathbf{q}%
_{2},\mathbf{q}_{3})$ de l’espai vectorial $\mathbb{R}_{2}\left[ x\right] $%
, essent $\mathbf{p}_{1}=2+x^{2},\mathbf{p}_{2}=1+2x+x^{2},\mathbf{p}%
_{3}=2x+x^{2},\mathbf{p}_{4}=1+x,\mathbf{q}_{1}=1,\mathbf{q}_{2}=x,\mathbf{q}%
_{3}=x^{2}$. Demostreu que són equivalents.

\end{exercici}
\begin{solucio}
Per demostrar-ho veurem que els vectors de cada sistema són combinació lineal dels de l'altre. En efecte, és fàcil trobar les següents
relacions%
\begin{align*}
\mathbf{q}_{1} &=\mathbf{p}_{2}-\mathbf{p}_{3} \\
\mathbf{q}_{2} &=-\mathbf{p}_{2}+\mathbf{p}_{3}+\mathbf{p}_{4} \\
\mathbf{q}_{3} &=\mathbf{p}_{1}-2\mathbf{p}_{2}+2\mathbf{p}_{3}
\end{align*}%
i%
\begin{align*}
\mathbf{p}_{1} &=2\mathbf{q}_{1}+\mathbf{q}_{3} \\
\mathbf{p}_{2} &=\mathbf{q}_{1}+2\mathbf{q}_{2}+\mathbf{q}_{3} \\
\mathbf{p}_{3} &=2\mathbf{q}_{2}+\mathbf{q}_{3} \\
\mathbf{p}_{4} &=\mathbf{q}_{1}+\mathbf{q}_{2}
\end{align*}%
Per tant, $\mathcal{S}_{1}\sim \mathcal{S}_{2}$.
\end{solucio}

\section{Dependència i independència lineal}

\begin{exercici}
Quines de les següents famílies finites d’elements
de $\mathcal{F}(\mathbb{R},\mathbb{R})$ són lliures? (a) $(f_{1},f_{2},f_{3})$%
, definides de la manera següent: $f_{1}(x)=\sin x,f_{2}(x)=\cos x,f_{3}(x)=1$; (b) $%
(g_{1},g_{2})$, definides per $g_{1}(x)=e^{x},g_{2}(x)=e^{x+2}$; (c) $%
(h_{1},h_{2},h_{3})$, definides com $h_{1}(x)=2,h_{2}(x)=x+2,h_{3}(x)=x^{2}$%
; (d) $(j_{1},j_{2},j_{3})$, $j_{1}(x)=0,j_{2}(x)=1,j_{3}(x)=x+1$.

\end{exercici}
\begin{solucio}
(a) Considerem%
\begin{equation*}
\alpha _{1}f_{1}+\alpha _{2}f_{2}+\alpha _{3}f_{3}=\mathbf{0}
\end{equation*}%
Aleshores, per a tot $x\in \mathbb{R}$ es compleix%
\begin{equation*}
\alpha _{1}\sin x+\alpha _{2}\cos x+\alpha _{3}1=0
\end{equation*}%
En particular, per a $x=0,x=\pi /2,x=\pi $, es té%
\begin{equation*}
\left.
\begin{array}{r}
\alpha _{2}+\alpha _{3}=0 \\
\alpha _{1}+\alpha _{3}=0 \\
-\alpha _{2}+\alpha _{3}=0%
\end{array}%
\right\}
\end{equation*}%
la solució del qual és $\alpha _{3}=0,\alpha _{2}=0,\alpha _{1}=0$. Per tant,
la família és lliure.

(b) Considerem%
\begin{equation*}
\alpha _{1}g_{1}+\alpha _{2}g_{2}=\mathbf{0}
\end{equation*}%
Aleshores, per a tot $x\in \mathbb{R}$ es compleix%
\begin{equation*}
\alpha _{1}e^{x}+\alpha _{2}e^{x+2}=0
\end{equation*}%
és a dir,%
\begin{equation*}
(\alpha _{1}+\alpha _{2}e^{2})e^{x}=0
\end{equation*}%
De on%
\begin{equation*}
\alpha _{1}=-e^{2}\alpha _{2}
\end{equation*}%
Prenent $\alpha _{2}=1$ i $\alpha _{1}=-e^{2}$, es té%
\begin{equation*}
-e^{2}g_{1}+g_{2}=0
\end{equation*}%
Per tant, la família és lligada.

(c) Considerem%
\begin{equation*}
\alpha _{1}h_{1}+\alpha _{2}h_{2}+\alpha _{3}h_{3}=\mathbf{0}
\end{equation*}%
Aleshores, per a tot $x\in \mathbb{R}$, es compleix%
\begin{equation*}
\alpha _{1}2+\alpha _{2}(x+2)+\alpha _{3}x^{2}=0
\end{equation*}%
és a dir,%
\begin{equation*}
2(\alpha _{1}+\alpha _{2})+\alpha _{2}x+\alpha _{3}x^{2}=0
\end{equation*}%
De on s’obté%
\begin{equation*}
\left.
\begin{array}{r}
2(\alpha _{1}+\alpha _{2})=0 \\
\alpha _{2}=0 \\
\alpha _{3}=0%
\end{array}%
\right\}
\end{equation*}%
la solució del qual és $\alpha _{1}=\alpha _{2}=\alpha _{3}=0$. Per tant, la
família és lliure.

(d) La família és lligada ja que conté el vector zero.
\end{solucio}

\begin{exercici}
Trobeu els valors de $a$ i $b$ perquè els vectors $%
(3,2,a,5),(2,-3,5,a),(0,13,b,7)$ de $\mathbb{R}^{4}$ siguin linealment
dependents.

\end{exercici}
\begin{solucio}
Considerem%
\begin{equation*}
\lambda _{1}(3,2,a,5)+\lambda _{2}(2,-3,5,a)+\lambda _{3}(0,13,b,7)=(0,0,0,0)
\end{equation*}%
és a dir,%
\begin{equation*}
(3\lambda _{1}+2\lambda _{2},2\lambda _{1}-3\lambda _{2}+13\lambda
_{3},a\lambda _{1}+5\lambda _{2}+b\lambda _{3},5\lambda _{1}+a\lambda
_{2}+7\lambda _{3})=(0,0,0,0)
\end{equation*}%
Per tant,%
\begin{equation*}
\begin{array}{r}
3\lambda _{1}+2\lambda _{2}=0 \\
2\lambda _{1}-3\lambda _{2}+13\lambda _{3}=0 \\
a\lambda _{1}+5\lambda _{2}+b\lambda _{3}=0 \\
5\lambda _{1}+a\lambda _{2}+7\lambda _{3}=0%
\end{array}%
\end{equation*}%
De les dues primeres equacions s’obté%
\begin{equation*}
\lambda _{2}=-\frac{3}{2}\lambda _{1}\text{ \ i \ }\lambda _{3}=-\frac{1}{2}%
\lambda _{1}
\end{equation*}%
Substituint aquests resultats en les dues següents i suposant que $%
\lambda _{1}\neq 0$ (ja que en cas contrari serien linealment
independents) s’obté%
\begin{equation*}
\left.
\begin{array}{r}
2a-b=15 \\
-3a=-3%
\end{array}%
\right\}
\end{equation*}%
d’on resulta%
\begin{equation*}
a=1\text{ \ i \ }b=-13
\end{equation*}
\end{solucio}

\begin{exercici}
En l’espai vectorial dels nombres reals $\mathbb{R}$ sobre el
cos dels nombres racionals $\mathbb{Q}$, proveu que els vectors $%
1,\sqrt{2},\sqrt{5}$ són linealment independents.

\end{exercici}
\begin{solucio}
Considerem%
\begin{equation*}
\alpha _{1}1+\alpha _{2}\sqrt{2}+\alpha _{3}\sqrt{5}=0
\end{equation*}%
amb $\alpha _{1},\alpha _{2},\alpha _{3}\in \mathbb{Q}$. Aïllant $%
\alpha _{2}\sqrt{2}$ i elevant al quadrat, s'obté%
\begin{equation*}
\alpha _{2}\sqrt{2}=-\alpha _{1}-\alpha _{3}\sqrt{5}\qquad \Longrightarrow
\qquad 2\alpha _{2}^{2}=\alpha _{1}^{2}+5\alpha _{3}^{2}+2\alpha _{1}\alpha
_{3}\sqrt{5}
\end{equation*}%
Si fos $\alpha _{1}\alpha _{3}\neq 0$, aleshores%
\begin{equation*}
\sqrt{5}=\frac{2\alpha _{2}^{2}-\alpha _{1}^{2}-5\alpha _{3}^{2}}{2\alpha
_{1}\alpha _{3}}
\end{equation*}%
seria un nombre racional, cosa que no és possible. Per tant, $\alpha _{1}=0$
o $\alpha _{3}=0$. Considerem els dos casos:

Si $\alpha _{3}=0$, la relació és redueix a $\alpha _{1}+\alpha _{2}\sqrt{2}%
=0$. Si fos $\alpha _{2}\neq 0$, aleshores $\sqrt{2}=-\alpha _{1}/\alpha
_{2} $ seria racional, cosa que no és possible. Per tant, $\alpha _{2}=0$ i,
en conseqüència, $\alpha _{1}=0$.

Si $\alpha _{1}=0$, la relació és redueix a $\alpha _{2}\sqrt{2}+\alpha _{3}%
\sqrt{5}=0$. Si fos $\alpha _{3}\neq 0$, elevant al quadrat s'obtindria $%
2\alpha _{2}^{2}=5\alpha _{3}^{2}$, és a dir, $\left( \alpha _{2}/\alpha
_{3}\right) ^{2}=5/2$, cosa que no és possible per a cap nombre racional,
perquè l'equació $2p^{2}=5q^{2}$ no té solucions enteres no nul\textperiodcentered les
(unicitat de la factorització en nombres primers). Per tant, $\alpha _{3}=0$
i, pel cas anterior, $\alpha _{2}=\alpha _{1}=0$.

En conseqüència, els vectors són linealment independents.
\end{solucio}

\section{Bases d’un espai vectorial de dimensió finita}

\begin{exercici}
Demostreu que els vectors $\mathbf{u}_{1}=(1,1,1),\mathbf{u}_{2}=(1,1,2)$ i
$\mathbf{u}_{3}=(1,2,3)$ formen una base de $\mathbb{R}^{3}$, i trobeu
les coordenades del vector $\mathbf{v}=(6,9,14)$ en aquesta base.

\end{exercici}
\begin{solucio}
Sabem que la dimensió de l’espai vectorial és tres. Per tant, els
vectors formaran una basi són linealment independents. Per veure-ho,
considerem%
\begin{equation*}
\lambda _{1}(1,1,1)+\lambda _{2}(1,1,2)+\lambda _{3}(1,2,3)=(0,0,0)
\end{equation*}%
és a dir,%
\begin{equation*}
(\lambda _{1}+\lambda _{2}+\lambda _{3},\lambda _{1}+\lambda _{2}+2\lambda
_{3},\lambda _{1}+2\lambda _{2}+3\lambda _{3})=(0,0,0)
\end{equation*}%
Per tant, s’obté el sistema següent%
\begin{equation*}
\left.
\begin{array}{r}
\lambda _{1}+\lambda _{2}+\lambda _{3}=0 \\
\lambda _{1}+\lambda _{2}+2\lambda _{3}=0 \\
\lambda _{1}+2\lambda _{2}+3\lambda _{3}=0%
\end{array}%
\right\}
\end{equation*}%
la solució del qual és $\lambda _{1}=\lambda _{2}=\lambda _{3}=0$. Per tant,
els vectors són linealment independents i, en conseqüència, formen base.
Las coordenades $(x,y,z)$ de $\mathbf{v}$ en aquesta base compleixen%
\begin{equation*}
x(1,1,1)+i(1,1,2)+z(1,2,3)=(6,9,14)
\end{equation*}%
és a dir,%
\begin{equation*}
\left.
\begin{array}{r}
x+y+z=6 \\
x+i+2z=9 \\
x+2y+3z=14%
\end{array}%
\right\}
\end{equation*}%
la solució del qual és $x=1,i=2$ i $z=3$. Per consegüent, les coordenades de
$\mathbf{v}$ són $(1,2,3)$ en la base en cuestión.
\end{solucio}

\begin{exercici}
En l’espai vectorial $\mathcal{F}(\mathbb{R},\mathbb{R})$, es consideren
les funcions $f_{1},f_{2},f_{3},f_{4},f_{5}$ definides de la manera següent: $%
f_{1}(x)=1,f_{2}(x)=sin^{2}x,f_{3}(x)=cos^{2}x,f_{4}(x)=sin2x,f_{5}(x)=cos2x$%
. Es demana: (a) estudiar-ne la dependència o independència lineal; (b)
determineu la dimensió i una base del subespai $F%
=\left\langle f_{1},f_{2},f_{3},f_{4},f_{5}\right\rangle $; (c) calculeu les
coordenades de $f$ i $g$ amb respecte de la base trobada en (b), essent $%
f(x)=cos2x+sin2x$ i $g(x)=cosx$.

\end{exercici}
\begin{solucio}
(a) Els vectors són linealment dependents, ja que es té%
\begin{align*}
f_{3} &=\frac{1}{2}f_{1}+\frac{1}{2}f_{5} \\
f_{2} &=\frac{1}{2}f_{1}-\frac{1}{2}f_{5}
\end{align*}%
ja que $cos^{2}x=\frac{1}{2}(1+cos2x)$ i $sin^{2}x=\frac{1}{2}(1-cos2x)$.

(b) Per les relacions anteriors i, aplicant transformacions elementals,
es té%
\begin{equation*}
\left( f_{1},f_{2},f_{3},f_{4},f_{5}\right) \sim \left(
f_{1},f_{4},f_{5}\right)
\end{equation*}%
Per tant,%
\begin{equation*}
F=\left\langle f_{1},f_{2},f_{3},f_{4},f_{5}\right\rangle
=\left\langle f_{1},f_{4},f_{5}\right\rangle
\end{equation*}%
Provarem que $f_{1},f_{4},f_{5}$ són linealment independents.
Considerem%
\begin{equation*}
\lambda _{1}f_{1}+\lambda _{2}f_{4}+\lambda _{3}f_{5}=\mathbf{0}
\end{equation*}%
és a dir, per a tot $x\in \mathbb{R}$ es compleix%
\begin{equation*}
\lambda _{1}1+\lambda _{2}sin2x+\lambda _{3}cos2x=\mathbf{0}
\end{equation*}%
Prenent $x=0,x=\pi /4,x=\pi /2$ es té el sistema següent%
\begin{equation*}
\left.
\begin{array}{r}
\lambda _{1}+\lambda _{3}=0 \\
\lambda _{1}+\lambda _{2}=0 \\
\lambda _{1}-\lambda _{3}=0%
\end{array}%
\right\}
\end{equation*}%
la solució del qual és $\lambda _{1}=\lambda _{2}=\lambda _{3}=0$. En
conseqüència, $dim\,F=3$ i $(f_{1},f_{4},f_{5})$ és una base de $%
F$.

(c) És clar que
\begin{equation*}
f=0f_{1}+1f_{4}+1f_{5}
\end{equation*}%
Per tant, les coordenades de $f$ en la base considerada són $(0,1,1)$. D’altra banda, es compleix $g\notin F$ ja que, en cas contrari, tindríem per a tot $x\in \mathbb{R}$%
\begin{equation*}
cosx=\alpha +\beta sin2x+\gamma cos2x
\end{equation*}%
Prenent $x=0$ i $x=\pi $, s’obtindria%
\begin{align*}
1 &=\alpha +\gamma \\
-1 &=\alpha +\gamma
\end{align*}%
la qual cosa és una contradicció. Per consegüent, no té sentit parlar
de coordenades en $F$.
\end{solucio}

\section{Dimensió d’un subespai vectorial}

\begin{exercici}
Considerem l’espai vectorial $\mathbb{R}_{n}\left[ x\right] $ dels
polinomis amb coeficients reals de grau menor o igual que $n$. Se pide:
(a) demostreu que el polinomi $x^{n}$ i les seves $n$ primeres derivades formen
una base de $\mathbb{R}_{n}\left[ x\right] $; (b) estudiar si els vectors $%
p_{1}(x)=1+3x+5x^{2},p_{2}(x)=-1+2x^{2},p_{3}(x)=3+3x+x^{2}$ són o no
linealment independents; (c) siguin $r_{1}(x)=1+x^{2},r_{2}(x)=1-x^{2}$ i $%
S_{1}=\left\langle r_{1}(x),r_{2}(x)\right\rangle $, estudiar si $%
p(x)$ i $r(x)$ pertanyen a $S_{1}$, essent $p(x)=1+5x^{2},r(x)=1+x$%
; (d) sigui $S_{2}=\left\langle
p_{1}(x),p_{2}(x),p_{3}(x)\right\rangle $, calculeu $S_{1}\cap
S_{2}$ i $S_{1}+S_{2}$.

\end{exercici}
\begin{solucio}
(a) Se compleix:%
\begin{align*}
f_{0}(x) &=x^{n} \\
f_{1}(x) &=nx^{n-1} \\
f_{2}(x) &=n(n-1)x^{n-2} \\
&&\vdots \\
f_{n-1}(x) &=n(n-1)(n-2)\cdots 2x \\
f_{n}(x) &=n!
\end{align*}%
Hay que proveu que $(f_{0}(x),f_{1}(x),...,f_{n}(x))$ és una base de $%
\mathbb{R}_{n}\left[ x\right] $. Qualsevol polinomi $p(x)$ de $\mathbb{R}%
_{n}\left[ x\right] $ és de la forma%
\begin{equation*}
p(x)=a_{0}+a_{1}x+\cdots +a_{n-1}x^{n-1}+a_{n}x^{n}
\end{equation*}%
expressant-se també de la manera següent:%
\begin{equation*}
p(x)=\frac{a_{o}}{n!}f_{n}(x)+\frac{a_{1}}{n(n-1)\cdots 2}f_{n-1}(x)+\cdots +%
\frac{a_{n-1}}{n}f_{1}(x)+a_{n}f_{0}(x)
\end{equation*}%
és a dir, $p(x)$ és combinació lineal de $%
f_{0}(x),f_{1}(x),...,f_{n}(x) $. D’altra banda, sabem que $dim\,\mathbb{R%
}_{n}\left[ x\right] =n+1$. Per tant, $(f_{0}(x),f_{1}(x),...,f_{n}(x))$ és
una base.

(b) Observeu que $p_{3}(x)=p_{1}(x)-p_{2}(x)$. Per tant, són
linealment dependents.

(c) Perquè $p(x)\in \left\langle r_{1}(x),r_{2}(x)\right\rangle $ han de
existir $\lambda _{1},\lambda _{2}\in \mathbb{R}$ tals que%
\begin{equation*}
p(x)=\lambda _{1}r_{1}(x)+\lambda _{2}r_{2}(x)
\end{equation*}%
és a dir,%
\begin{equation*}
1+5x^{2}=\lambda _{1}(1+x^{2})+\lambda _{2}(1-x^{2})
\end{equation*}%
D’aquí s’obté el sistema%
\begin{equation*}
\left.
\begin{array}{r}
\lambda _{1}+\lambda _{2}=1 \\
\lambda _{1}-\lambda _{2}=5%
\end{array}%
\right\}
\end{equation*}%
la solució del qual és $\lambda _{2}=-2,\lambda _{1}=3$. Per tant, $p(x)\in
S_{1}$. De la mateixa manera, $r(x)\in \left\langle
r_{1}(x),r_{2}(x)\right\rangle $ ha de complirse%
\begin{equation*}
r(x)=\mu _{1}r_{1}(x)+\mu _{2}r_{2}(x)
\end{equation*}%
és a dir,%
\begin{equation*}
1+x=\mu _{1}(1+x^{2})+\mu _{2}(1-x^{2})
\end{equation*}%
o bé,%
\begin{equation*}
1+x=\mu _{1}+\mu _{2}+0x+(\mu _{1}-\mu _{2})x^{2}
\end{equation*}%
el que no és possible, ja que $r(x)$ té coeficient $1$ en $x$ i, en
canvi, $\mu _{1}r_{1}(x)+\mu _{2}r_{2}(x)$ té coeficient $0$ en $x$.
Per tant, $r(x)\notin S_{1}$.

(d) Per (b) es té $S_{2}=\left\langle
p_{1}(x),p_{2}(x)\right\rangle $ i $dim\,S_{2}=2$. D'altra banda, $%
(r_{1}(x),r_{2}(x))$ és una base de $S_{1}$ i $dim\,S%
_{1}=2 $. Per determinar $p(x)$ de manera que $p(x)\in S_{1}\cap
S_{2}$ ha de complirse%
\begin{equation*}
p(x)=\alpha p_{1}(x)+\beta p_{2}(x)\text{ \ \ \ i \ \ \ }p(x)=\lambda
r_{1}(x)+\mu r_{2}(x)
\end{equation*}%
és a dir,%
\begin{equation*}
\alpha (1+3x+5x^{2})+\beta (-1+2x^{2})=\lambda (1+x^{2})+\mu (1-x^{2})
\end{equation*}%
o bé,%
\begin{equation*}
\alpha -\beta +3\alpha x+(5\alpha +2\beta )x^{2}=\lambda +\mu +(\lambda -\mu
)x^{2}
\end{equation*}%
Per tant, tenim el següent sistema%
\begin{equation*}
\left.
\begin{array}{c}
\alpha -\beta =\lambda +\mu \\
3\alpha =0 \\
5\alpha +2\beta =\lambda -\mu%
\end{array}%
\right\}
\end{equation*}%
la solució del qual és $\alpha =0,\lambda =\frac{1}{2}\beta ,\mu =-\frac{3}{2}%
\beta $. Prenent $\beta =1$, un vector de la intersecció és $-1+2x^{2}$%
. Per consegüent, $S_{1}\cap S_{2}=\left\langle
-1+2x^{2}\right\rangle $ i $dim\,(S_{1}\cap S_{2})=1$.
Aleshores%
\begin{equation*}
dim(S_{1}+S_{2})=dim\,S_{1}+dim\,S%
_{2}-dim(S_{1}\cap S_{2})=2+2-1=3
\end{equation*}%
És clar que $S_{1}+S_{2}=\left\langle
p_{1}(x),p_{2}(x),r_{1}(x),r_{2}(x)\right\rangle $. Ara bé, només
tres d’aquests vectors són linealment independents. Per trobar-los,
considerem%
\begin{equation*}
\alpha (1+3x+5x^{2})+\beta (-1+2x^{2})+\gamma (1+x^{2})+\delta (1-x^{2})=%
\mathbf{0}
\end{equation*}%
d’on s’obté%
\begin{equation*}
(\alpha -\beta +\gamma +\delta )+3\alpha x+(5\alpha +2\beta +\gamma -\delta
)x^{2}=\mathbf{0}
\end{equation*}%
Per tant, ha de complirse%
\begin{equation*}
\left.
\begin{array}{r}
\alpha -\beta +\gamma +\delta =0 \\
3\alpha =0 \\
5\alpha +2\beta +\gamma -\delta =0%
\end{array}%
\right\}
\end{equation*}%
D’aquí s’obté, $\alpha =0$, $\beta =-2\gamma $ i $\delta =-3\gamma $%
. Prenent, per exemple, $\gamma =1$, aleshores $\beta =-2$ i $\delta =-3$ i,
en conseqüència, tenim $-2p_{2}(x)+r_{1}(x)-3r_{2}(x)=0$, és a dir,%
\begin{equation*}
r_{2}(x)=\frac{1}{3}r_{1}(x)-\frac{2}{3}p_{2}(x)
\end{equation*}%
és a dir, $r_{2}(x)$ depèn linealment dels altres. Per tant, $%
(p_{1}(x),p_{2}(x),r_{1}(x))$ és una base de $S_{1}+S_{2}$.
\end{solucio}

\section{Espai vectorial quocient}

\begin{exercici}
Sigui $E$ un espai vectorial de dimensió 3 i $%
(e_{1},e_{2},e_{3})$ una base de $E$. Sigui $F$ el
subespai generat per $(v_{1},v_{2})$ essent $v_{1}=e_{1}-e_{2}$ i $%
v_{2}=e_{1}+e_{2}$. Trobeu una base de l’espai vectorial quocient $E/F$.

\end{exercici}
\begin{solucio}
És clar que el sistema $(v_{1},v_{2})$ és lliure. Per tant, $dim\,F=2$. En
conseqüència,%
\begin{equation*}
dim\text{\thinspace }E/F=dim\,E-dim\,F%
=3-2=1
\end{equation*}%
Per trobeu una base de $E/F$ n’hi ha prou de trobar un vector no
nul de $E/F$, és a dir, una classe no nul·la. Si $\left[ v%
\right] \neq \left[ 0\right] $, aleshores $v\notin F$. Per tant, $v$
pot ser qualsevol vector que sigui linealment independent amb $v_{1}$ i $%
v_{2}$; per exemple, podem tomar $e_{3}$. Per consegüent, $(e_{3}+%
F)$ és una base de $E/F$.
\end{solucio}

\begin{exercici}
Obtener una base de l’espai vectorial quocient $\mathbb{R}^{4}$ mòdul
el subespai
\begin{equation*}
F=\left\langle (1,0,1,1),(1,2,1,1),(2,2,2,2)\right\rangle
\end{equation*}

\end{exercici}
\begin{solucio}
Observem que%
\begin{equation*}
(2,2,2,2)=(1,0,1,1)+(1,2,1,1)
\end{equation*}%
Per tant, $dim\,F=2$. D’aquí,
\begin{equation*}
dim\,\mathbb{R}^{4}/F=4-2=2
\end{equation*}%
Per obtenir una base de $\mathbb{R}^{4}/F$ n’hi ha prou de trobar una
base del subespai suplementari de $F$ respecte de $\mathbb{R}%
^{4} $. Obtindrem aquesta base completant la base de $F$ a una base
de $\mathbb{R}^{4}$. Sigui $(e_{1},e_{2},e_{3},e_{4})$ la base canònica de
$\mathbb{R}^{4}$ i $v_{1}=(1,0,1,1),v_{2}=(1,2,1,1)$. Observem que%
\begin{equation*}
e_{1}=v_{1}-e_{3}-e_{4}
\end{equation*}%
Per tant, $(v_{1},e_{2},e_{3},e_{4})$ és una base de $\mathbb{R}^{4}$.
Observem també que%
\begin{equation*}
e_{2}=-\frac{1}{2}v_{1}+\frac{1}{2}v_{2}
\end{equation*}%
Per tant, $(v_{1},v_{2},e_{3},e_{4})$ és una base de $\mathbb{R}^{4}$. En
conseqüència, $(e_{3},e_{4})$ és una base del subespai suplementari de $%
F$. Per consegüent, $(e_{3}+F,e_{4}+F)$ és una
base de $\mathbb{R}^{4}/F$.
\end{solucio}

\section{Rang d’un sistema de vectors}

\begin{exercici}
Determineu, segons els valors de $\alpha $, el rang del sistema $%
\mathcal{S}=(\mathbf{v}_{1},\mathbf{v}_{2},\mathbf{v}_{3})$ de vectors de $%
\mathbb{R}^{3}$, essent $\mathbf{v}_{1}=(\alpha ,-1,-1),\mathbf{v}%
_{2}=(-1,-\alpha ,-1),\mathbf{v}_{3}=(-1,-1,\alpha )$.

\end{exercici}
\begin{solucio}
Prenent la base canònica de $\mathbb{R}^{3}$, escrivim les coordenades
d’aquests vectors per files (de la mateixa manera es faria per columnes) com
sigue:%
\begin{equation*}
\begin{array}{rrrr}
\mathbf{v}_{1}: & \alpha & -1 & -1 \\
\mathbf{v}_{2}: & -1 & -\alpha & -1 \\
\mathbf{v}_{3}: & -1 & -1 & \alpha%
\end{array}%
\end{equation*}%
Aplicant T1 s’obté%
\begin{equation*}
\begin{array}{rrrr}
\mathbf{v}_{3}: & -1 & -1 & \alpha \\
\mathbf{v}_{2}: & -1 & -\alpha & -1 \\
\mathbf{v}_{1}: & \alpha & -1 & -1%
\end{array}%
\end{equation*}%
aplicant 2 veces T3 s’obté%
\begin{equation*}
\begin{array}{rrrr}
\mathbf{v}_{3}: & -1 & -1 & \alpha \\
\mathbf{v}_{2}-\mathbf{v}_{3}: & 0 & -\alpha +1 & -1-\alpha \\
\mathbf{v}_{1}+\alpha \mathbf{v}_{3}: & 0 & -1-\alpha & -1+\alpha ^{2}%
\end{array}%
\end{equation*}%
o de manera equivalent%
\begin{equation*}
\begin{array}{rrrr}
\mathbf{v}_{3}: & -1 & -1 & \alpha \\
\mathbf{v}_{2}-\mathbf{v}_{3}: & 0 & 1-\alpha & -(1+\alpha ) \\
\mathbf{v}_{1}+\alpha \mathbf{v}_{3}: & 0 & -(\alpha +1) & (\alpha
+1)(\alpha -1)%
\end{array}%
\end{equation*}%
Aplicant T2, suposant que $\alpha +1\neq 0$, s’obté%
\begin{equation*}
\begin{array}{rrrr}
\mathbf{v}_{3}: & -1 & -1 & \alpha \\
\mathbf{v}_{2}-\mathbf{v}_{3}: & 0 & 1-\alpha & -(1+\alpha ) \\
\frac{\mathbf{v}_{1}+\alpha \mathbf{v}_{3}}{\alpha +1}: & 0 & -1 & \alpha -1%
\end{array}%
\end{equation*}%
Aplicant T1, es té%
\begin{equation*}
\begin{array}{rrrr}
\mathbf{v}_{3}: & -1 & -1 & \alpha \\
\frac{\mathbf{v}_{1}+\alpha \mathbf{v}_{3}}{\alpha +1}: & 0 & -1 & \alpha -1
\\
\mathbf{v}_{2}-\mathbf{v}_{3}: & 0 & 1-\alpha & -(1+\alpha )%
\end{array}%
\end{equation*}%
aplicant T3, es té%
\begin{equation*}
\begin{array}{rrrr}
\mathbf{v}_{3}: & -1 & -1 & \alpha \\
\frac{\mathbf{v}_{1}+\alpha \mathbf{v}_{3}}{\alpha +1}: & 0 & -1 & \alpha -1
\\
\mathbf{v}_{2}-\mathbf{v}_{3}+\frac{(1-\alpha )(\mathbf{v}_{1}+\alpha
\mathbf{v}_{3})}{\alpha +1}: & 0 & 0 & -(1+\alpha )-(\alpha -1)^{2}%
\end{array}%
\end{equation*}%
o de manera equivalent%
\begin{equation*}
\begin{array}{rrrr}
\mathbf{v}_{3}: & -1 & -1 & \alpha \\
\frac{\mathbf{v}_{1}+\alpha \mathbf{v}_{3}}{\alpha +1}: & 0 & -1 & \alpha -1
\\
\mathbf{v}_{2}-\mathbf{v}_{3}+\frac{(1-\alpha )(\mathbf{v}_{1}+\alpha
\mathbf{v}_{3})}{\alpha +1}: & 0 & 0 & -\alpha ^{2}+\alpha -2%
\end{array}%
\end{equation*}%
De on obtenim que si $\alpha \neq -1$, aleshores%
\begin{equation*}
(\mathbf{v}_{1},\mathbf{v}_{2},\mathbf{v}_{3})\sim \left( \mathbf{v}_{3},%
\frac{\mathbf{v}_{1}+\alpha \mathbf{v}_{3}}{\alpha +1},\mathbf{v}_{2}-%
\mathbf{v}_{3}+\frac{(1-\alpha )(\mathbf{v}_{1}+\alpha \mathbf{v}_{3})}{%
\alpha +1}\right)
\end{equation*}%
i, els vectors d’aquest darrer sistema són linealment independents,
ja que $-\alpha ^{2}+\alpha -2\neq 0$ per a tot $\alpha \in \mathbb{R}$. Per
tanto, si $\alpha \neq -1$ aleshores $rang\,\mathcal{S}=3$. Si fuese $\alpha
=-1$, aleshores tindríem%
\begin{equation*}
\begin{array}{rrrr}
\mathbf{v}_{3}: & -1 & -1 & -1 \\
\mathbf{v}_{2}-\mathbf{v}_{3}: & 0 & 2 & 0 \\
\mathbf{v}_{1}-\mathbf{v}_{3}: & 0 & 0 & 0%
\end{array}%
\end{equation*}%
Suprimint la tercera fila es té%
\begin{equation*}
\begin{array}{rrrr}
\mathbf{v}_{3}: & -1 & -1 & -1 \\
\mathbf{v}_{2}-\mathbf{v}_{3}: & 0 & 2 & 0%
\end{array}%
\end{equation*}%
Per tant, si $\alpha =-1$
\begin{equation*}
(\mathbf{v}_{1},\mathbf{v}_{2},\mathbf{v}_{3})\sim (\mathbf{v}_{3},\mathbf{v}%
_{2}-\mathbf{v}_{3})
\end{equation*}%
i es té $rang\,\mathcal{S}=2$. A més, la relació de dependència
en aquest cas és%
\begin{equation*}
\mathbf{v}_{1}-\mathbf{v}_{3}=\mathbf{0}
\end{equation*}
\end{solucio}

\begin{exercici}
Determineu el rang del sistema $\mathcal{S}$ de $\mathbb{R}^{4}$ i
si escau, les relacions de dependència entre els vectors del sistema,
essent $\mathcal{S}=(\mathbf{v}_{1},\mathbf{v}_{2},\mathbf{v}_{3},\mathbf{v}%
_{4})$ amb $\mathbf{v}_{1}=(1,1,1,1),\mathbf{v}_{2}=(0,1,2,-1),\mathbf{v}%
_{3}=(1,0,-2,3),\mathbf{v}_{4}=(2,1,0,-1)$.

\end{exercici}
\begin{solucio}
Prenent la base canònica de $\mathbb{R}^{4}$, escrivim les coordenades
d’aquests vectors per columnes (de la mateixa manera es faria per files) com
sigue:%
\begin{equation*}
\begin{array}{rrrr}
\mathbf{v}_{1} & \mathbf{v}_{2} & \mathbf{v}_{3} & \mathbf{v}_{4} \\
1 & 0 & 1 & 2 \\
1 & 1 & 0 & 1 \\
1 & 2 & -2 & 0 \\
1 & -1 & 3 & -1%
\end{array}%
\end{equation*}%
Aplicant dues vegades T3, es té%
\begin{equation*}
\begin{array}{rrrr}
\mathbf{v}_{1} & \mathbf{v}_{2} & \mathbf{v}_{3}-\mathbf{v}_{1} & \mathbf{v}%
_{4}-2\mathbf{v}_{1} \\
1 & 0 & 0 & 0 \\
1 & 1 & -1 & -1 \\
1 & 2 & -3 & -2 \\
1 & -1 & 2 & -3%
\end{array}%
\end{equation*}%
Aplicant de nou dues vegades T3, es té%
\begin{equation*}
\begin{array}{rrrr}
\mathbf{v}_{1} & \mathbf{v}_{2} & \mathbf{v}_{3}-\mathbf{v}_{1}+\mathbf{v}%
_{2} & \mathbf{v}_{4}-2\mathbf{v}_{1}+\mathbf{v}_{2} \\
1 & 0 & 0 & 0 \\
1 & 1 & 0 & 0 \\
1 & 2 & -1 & 0 \\
1 & -1 & 1 & -4%
\end{array}%
\end{equation*}%
Per tant,%
\begin{equation*}
rang(\mathbf{v}_{1},\mathbf{v}_{2},\mathbf{v}_{3},\mathbf{v}_{4})=rang(%
\mathbf{v}_{1},\mathbf{v}_{2},-\mathbf{v}_{1}+\mathbf{v}_{2}+\mathbf{v}%
_{3},-2\mathbf{v}_{1}+\mathbf{v}_{2}+\mathbf{v}_{4})=4
\end{equation*}%
ja que els vectors del sistema $(\mathbf{v}_{1},\mathbf{v}_{2},-\mathbf{v}%
_{1}+\mathbf{v}_{2}+\mathbf{v}_{3},-2\mathbf{v}_{1}+\mathbf{v}_{2}+\mathbf{v}%
_{4})$ són linealment independents.
\end{solucio}

\section{Producte d’espais vectorials. Isomorfismes d’espais
vectorials}

\begin{exercici}
Siguin els espais vectorials $E$ i $F$ amb $E=%
F=\mathbb{R}^{2}$. Calculeu una base de l’espai vectorial producte
$E\times F$ asociada a la base natural de $\mathbb{R}^{2}$.

\end{exercici}
\begin{solucio}
Sabem que l’espai vectorial $E\times F$ és de dimensió 4, essent una base els vectors $\mathbf{v}_{1}=(\mathbf{e}_{1},%
\mathbf{0}),\mathbf{v}_{2}=(\mathbf{e}_{2},\mathbf{0}),\mathbf{v}_{3}=(%
\mathbf{0},\mathbf{e}_{1})$ i $\mathbf{v}_{4}=(\mathbf{0},\mathbf{e}_{2})$.

Sabem que
\begin{equation*}
\mathbb{R}^{2}\times \mathbb{R}^{2}\cong \mathbb{R}^{2}\oplus \mathbb{R}%
^{2}\cong (\mathbb{R}\oplus \mathbb{R})\oplus (\mathbb{R}\oplus \mathbb{R}%
)\cong \mathbb{R}\oplus \mathbb{R}\oplus \mathbb{R}\oplus \mathbb{R}\cong
\mathbb{R}^{4}
\end{equation*}%
D’aquí, podem identificar el vector $\mathbf{v}_{1}=(\mathbf{e}_{1},%
\mathbf{0})$ de $\mathbb{R}^{2}\times \mathbb{R}^{2}$ amb el vector $%
(1,0,0,0)$ de $\mathbb{R}^{4}$. D’aquesta manera, tenim
\begin{equation*}
((1,0,0,0),(0,1,0,0),(0,0,1,0),(0,0,0,1))
\end{equation*}%
és una base de l’espai producte.
\end{solucio}

\section{Canvi de coordenades}

\begin{exercici}
Calculeu la matriu de canvi de la base $\mathcal{B}%
_{1}=((1,1,0),(-1,0,2),(0,2,5))$ a la base $\mathcal{B}%
_{2}=((0,1,1),(1,1,1),(3,1,0))$ de $\mathbb{R}^{3}$.

\end{exercici}
\begin{solucio}
Sigui $\mathcal{B}_{0}$ la base natural de $\mathbb{R}^{3}$. Sabem que la
matriu de pas de $\mathcal{B}_{1}$ a $\mathcal{B}_{0}$ és%
\begin{equation*}
P=\left(
\begin{array}{ccc}
1 & -1 & 0 \\
1 & 0 & 2 \\
0 & 2 & 5%
\end{array}%
\right)
\end{equation*}%
i la matriu de pas de $\mathcal{B}_{2}$ a $\mathcal{B}_{0}$ és%
\begin{equation*}
Q=\left(
\begin{array}{ccc}
0 & 1 & 3 \\
1 & 1 & 1 \\
1 & 1 & 0%
\end{array}%
\right)
\end{equation*}%
Si $(x_{1},y_{1},z_{1})$ són les coordenades d’un vector en $\mathcal{B}%
_{1} $ i $(x_{2},y_{2},z_{2})$ són les coordenades del mateix vector en $%
\mathcal{B}_{2}$. Aleshores, es compleix%
\begin{equation*}
\left(
\begin{array}{ccc}
1 & -1 & 0 \\
1 & 0 & 2 \\
0 & 2 & 5%
\end{array}%
\right) \left(
\begin{array}{c}
x_{1} \\
y_{1} \\
z_{1}%
\end{array}%
\right) =\left(
\begin{array}{ccc}
0 & 1 & 3 \\
1 & 1 & 1 \\
1 & 1 & 0%
\end{array}%
\right) \left(
\begin{array}{c}
x_{2} \\
y_{2} \\
z_{2}%
\end{array}%
\right)
\end{equation*}%
com que les matrius de canvi són sempre invertibles, podem escriure%
\begin{equation*}
\left(
\begin{array}{ccc}
0 & 1 & 3 \\
1 & 1 & 1 \\
1 & 1 & 0%
\end{array}%
\right) ^{-1}\left(
\begin{array}{ccc}
1 & -1 & 0 \\
1 & 0 & 2 \\
0 & 2 & 5%
\end{array}%
\right) \left(
\begin{array}{c}
x_{1} \\
y_{1} \\
z_{1}%
\end{array}%
\right) =\left(
\begin{array}{c}
x_{2} \\
y_{2} \\
z_{2}%
\end{array}%
\right)
\end{equation*}%
és a dir%
\begin{equation*}
\left(
\begin{array}{ccc}
2 & -3 & -4 \\
-2 & 5 & 9 \\
1 & -2 & -3%
\end{array}%
\right) \left(
\begin{array}{c}
x_{1} \\
y_{1} \\
z_{1}%
\end{array}%
\right) =\left(
\begin{array}{c}
x_{2} \\
y_{2} \\
z_{2}%
\end{array}%
\right)
\end{equation*}%
Per tant, la matriu de pas de $\mathcal{B}_{1}$ a $\mathcal{B}_{2}$ és%
\begin{equation*}
Q^{-1}P=\left(
\begin{array}{ccc}
2 & -3 & -4 \\
-2 & 5 & 9 \\
1 & -2 & -3%
\end{array}%
\right)
\end{equation*}%
Observeu el següent diagrama per a recordar com es pot deduir
directamente la matriu en cuestión%
\begin{equation*}
\mathcal{B}_{2}\overset{Q^{-1}}{\longleftarrow }\mathcal{B}_{0}\overset{P}{%
\longleftarrow }\mathcal{B}_{1}
\end{equation*}%
No hi ha que olvidar que el diagrama per a qualsevol base $\mathcal{B}$ que no
sigui la natural és sempre el següent%
\begin{equation*}
\mathcal{B}_{0}\overset{P}{\longleftarrow }\mathcal{B}
\end{equation*}%
on la matriu $P$ té per columnes les coordenades dels vectors de
la base $\mathcal{B}$.
\end{solucio}